\section{Continuation III: a sufficient relative correlation estimate}
\label{sec:reduction}

Consider the photon-free model \(H^{\rm C}_{\Lambda,L}\) of
Lemma~\ref{lem:hf-energy}, on the exact neutral sector. Define
\begin{align*}
E^{\rm MB}(R,Z)
&=\inf\sigma\bigl(H^{\rm C}_{\Lambda,L}(R,Z)|_{\mathcal H_{Z_{\rm tot}}}\bigr),\\
E^{\rm HF}(R,Z)
&=\inf_{\substack{\gamma^2=\gamma\\
\rank\gamma=2M_L+Z_{\rm tot}}}\mathcal E_{R,Z}(\gamma).
\end{align*}
The Hartree--Fock infimum here is over exact-charge Slater determinants;
it is not the larger mean-charge relaxation. Set
\begin{equation}
c(R,Z)=E^{\rm HF}(R,Z)-E^{\rm MB}(R,Z)\ge0,
\end{equation}
and write \(c(0)\) for the case without nuclei.

\begin{proposition}[Exact identity and a sufficient comparison]
\label{prop:reduction}
Under the box-size and coupling hypotheses of Theorem~\ref{thm:mf},
\begin{align}
E^{\rm MB}(R,Z)-E^{\rm MB,\vac}
={}&[E^{\rm HF}(R,Z)-E^{\rm HF,\vac}]
-[c(R,Z)-c(0)]
\nonumber\\
\ge{}&-A^{\rm MF}_\Lambda K-[c(R,Z)-c(0)].
\label{eq:reduction}
\end{align}
Moreover \(c(0)=-E^{\rm MB,\vac}\in[0,B_\Lambda V]\).
Consequently the photon-free primary bound \eqref{eq:S} follows from
\begin{equation}
c(R,Z)\le c(0)+C_\Lambda K.
\label{eq:rel-corr}
\end{equation}
\end{proposition}
\begin{proof}
The equality is the definition of \(c\), and the lower bound uses
Theorem~\ref{thm:mf} restricted to the exact-charge Slater class.
Its vacuum energy is zero because this class contains \(P_-\).
The free-vacuum trial gives \(E^{\rm MB,\vac}\le0\), while
\eqref{eq:full-coulomb-bound}, with \(H_D\ge0\) and \(K=0\), gives
\(E^{\rm MB,\vac}\ge-B_\Lambda V\).
\end{proof}

\begin{remark}[Sufficiency is not equivalence]
Equation~\eqref{eq:reduction} is exact, but
\eqref{eq:rel-corr} is a sufficient estimate, not a necessary condition
deduced from \eqref{eq:S}. The Hartree--Fock nuclear energy can itself be
large and positive, for example near nuclear collisions. It may compensate
a larger correlation-energy difference. There is no proved assertion here
that either correlation energy has a specific asymptotic density.
Separate extensive error bounds do not give the cancellation in
\eqref{eq:rel-corr}.
\end{remark}

\subsection{What existing correlation and gap results do not yet supply}
The correlation estimates of Bach and Graf--Solovej
\cite{Bach1992,GS1994} are useful precedents for controlling the discrepancy
between many-body and Hartree--Fock energies in Coulomb systems.
Applying them to this filled-sea reference, and then controlling the
\emph{difference} of errors with and without nuclei, requires new estimates.
Their conclusions are not already \eqref{eq:rel-corr}.

De Roeck--Salmhofer \cite{DRS2019} prove persistence of a spectral gap
uniformly in the size of a fermionic lattice, for a gapped free Hamiltonian
and sufficiently weak interactions satisfying suitable decay hypotheses.
A purely fermionic lattice model with a fixed Yukawa screening length is
a plausible intermediate problem once those hypotheses and interaction
norms are verified. Removing screening is a separate limit: the relevant
decay constants need not remain bounded.

Even a uniform source-free gap does not establish either of the two
configuration-uniform Feshbach inputs in Section~\ref{sec:effective}.
Arbitrarily clustered sources can produce a strong external potential.
One still needs control of the relevant charged sectors and defects,
as well as an extensive bound on the virtual-pair correction.
Gap persistence alone does not give \eqref{eq:rel-corr}.

Global neutrality also does not provide a decomposition into localized
neutral objects with controlled dipole moments. If such a decomposition
were independently proved, dipole interactions could have an \(r^{-3}\)
tail. In three dimensions its absolute radial integral is logarithmic at
large distances. That is an infrared summability problem, and must not
be identified with the ultraviolet logarithm in charge renormalization.

\subsection{Gauge-covariant lattice regularization and gluing}
A gauge-covariant Hamiltonian lattice regularization is a reasonable
framework to investigate, but its kinematic and physical Hilbert spaces
must be distinguished. Site fermions and link variables give a kinematic
tensor product. The physical space is constrained by
\begin{equation}
\mathcal H_{\rm phys}
=\bigcap_x\ker\bigl[(\operatorname{div}E)_x-\rho_x-\nu_x\bigr].
\end{equation}
It does not in general factor as an ordinary tensor product across a
spatial boundary. Gluing must keep track of boundary electric-flux
sectors, or use an enlarged Hilbert space followed by the constraints
\cite{KS1975,Donnelly2014}. Gauge links and unbounded electric-field
operators also lie outside the purely fermionic setting of
\cite{DRS2019}.

With unitary link variables, a finite-range lattice hopping Hamiltonian
has a lower bound proportional to the number of sites; positive electric
and magnetic field energies can help preserve that coarse estimate.
This observation does not compare the source energy with the actual
interacting vacuum.

Exact lattice gauge covariance supplies the corresponding identities,
but does not make a spatially constant vector potential on a finite
torus gauge trivial. Holonomies and boundary conditions must be treated.
The field-dependent no-pair projection of Lieb--Loss \cite{LL2002}
changes the matter Hilbert space; it is not a regulator-equivalent
replacement for full pair-producing QED. The Pauli--Villars construction
of \cite{GHLS2013} concerns a Dirac vacuum with classical fields and small
external sources, not quantized-photon many-body stability.

\subsection{Scope of the ultraviolet question}
A Landau scale is a useful warning about the ultraviolet problem,
not a substitute for a constructive theorem. In the reduced, exchange-free
BDF setting, a charge-renormalization relation of the form
\begin{equation}
\alpha_{\rm phys}
=\frac{\alpha_{\rm bare}}{1+\alpha_{\rm bare} B_\Lambda},
\qquad
\alpha_{\rm bare}
=\frac{\alpha_{\rm phys}}{1-\alpha_{\rm phys}B_\Lambda},
\end{equation}
with \(B_\Lambda\sim(2/3\pi)\log(\Lambda/m)\), illustrates the restriction
\(\alpha_{\rm phys}B_\Lambda<1\) for positive bare coupling
\cite{HLS2005,GLS2009}. The associated exponential scale is model-specific and
based on the indicated asymptotics. It neither proves a full-QED Landau
pole nor establishes uniform stability up to \(m\exp(c/\alpha)\).
For this project, every proposed cutoff window should be stated explicitly
and its uniform estimates proved within a specified model.

\section{Status at the start of the further proof attempts}
\label{sec:status}

The table records the position after the audit of Sections 9--11.
Sections~\ref{uv-section}--\ref{sl-sec} subsequently settle three of
the proposed next steps in their specified models.
Section~\ref{next-status} gives the current conclusions.

\small
\noindent
\begin{tabular}{@{}>{\raggedright\arraybackslash}p{0.50\linewidth}
>{\raggedright\arraybackslash}p{0.44\linewidth}@{}}
\toprule
Statement & Audited status\\
\midrule
Finite-regulator self-adjointness and ground states
& Proved; Section~\ref{sec:lemmas}.\\[0.45em]
Massless baseline: lower bound and \(L^5\) vacuum trial
& Proved for the fixed momentum cutoff; Sections~\ref{sec:lower-proof}--\ref{sec:counterexamples}.\\[0.45em]
Free-sea static susceptibility and soft-momentum coefficient
& Proved; physical photon pole mass not identified.\\[0.45em]
Extensive bound with a positive bare photon-mass term
& Proved for the modified model; Theorem~\ref{thm:ct}.\\[0.45em]
Small-box failure under the two stated self-energy conventions
& Proved analytically; Remark~\ref{rem:madelung}.\\[0.45em]
Periodic Coulomb--Kato constant and Green-function remainder
& Explicit analytic bounds; Lemma~\ref{lem:pkato}.\\[0.45em]
Photon-free mean-field relative bound with point nuclei and pair coherence
& Proved for the stated coupling and box-size range; Theorem~\ref{thm:mf}.\\[0.45em]
Product mean field with a sufficiently large bare photon mass
& Proved with improved current estimate; Proposition~\ref{prop:mf-heavy}.\\[0.45em]
Uniform strict kinetic absorption at the one-loop mass tuning
& False: neutral Slater perturbations approach ratio one;
Remark~\ref{rem:mf-photon-open}.\\[0.45em]
Vacuum-relative theorem at that tuning
& Unresolved; failure of the preceding sufficient route does not refute it.\\[0.45em]
Many-body relative theorem from a relative correlation bound
& Historical reduction; the primary paper now proves the comparison
in its positive-effective-mass regime.\\[0.45em]
Cutoff-uniform stability, thermodynamic limit, binding, crystals
& Require additional arguments within each specified model.\\
\bottomrule
\end{tabular}
\normalsize

\subsection{Literature comparison and priority}
The finite-cutoff torus mean-field argument is presented as a direct
adaptation, with no claim of priority. In particular, the periodic
Coulomb comparison and Kato-type machinery already have substantial
precedent in Lemma~1 of Hainzl--Lewin--Solovej \cite{HLSol2007}.
Their periodic Coulomb kernel has an added constant chosen to be
nonnegative. They begin with symmetrized charge density rather than our
fixed free-vacuum Wick ordering, construct a self-consistent
Hartree--Fock vacuum, and derive the sea-relative BDF functional by
subtracting that vacuum. Their thermodynamic external-field statement
uses regular charge densities, and is not the arbitrary-\(K\),
point-source theorem asserted here.

Bach--Barbaroux--Helffer--Siedentop \cite{BBHS1998} include pair degrees of
freedom in generalized Hartree--Fock theory, point Coulomb nuclei, and
classical magnetic fields. Their splitting uses the positive spectral
space of an effective nuclear-field Dirac operator, and their magnetic
formulation also depends on the field. Their normal-ordering reference
therefore differs from the fixed \(P_-(D_0)\) reference used here.
This distinction prevents simply reading their theorem as our
fixed-vacuum comparison; it does not justify a claim that the present
stability result is new.

Frank--Lieb--Seiringer \cite{FLS2007} establish critical-charge
pseudorelativistic stability in magnetic fields. Their scalar
one-particle kinetic operator is not a quantized Dirac field with pairs.
These results should guide a cutoff-uniform investigation without being
substituted for the required theorem.

\subsection{Questions taken up in the following sections}
\begin{enumerate}
\item Could the photon-free mean-field bound be made cutoff independent
under subcritical point charges? Theorem~\ref{uv-theorem} disproves this
for the fixed free-Wick model, already with one unit nucleus.
\item Does a screened lattice model satisfy the relative energy and
correlation estimates? Theorem~\ref{sl-relative-theorem} proves both by
an exact-sector comparison, without a gap-persistence argument.
\item What is the vacuum at the one-loop static photon tuning?
Theorem~\ref{pv-instability-theorem} proves that it is not the free sea;
exchange gives an extensive variational improvement. Construction of
the actual reference vacuum remains a further problem.
\item If adopting gauge links, develop gluing with the Gauss-law
boundary-flux constraints. Add quantized photons only with explicit
control of their gapless sectors and of dressing errors.
\end{enumerate}
